% GENERATED FILE. Edit canonical lessons, then run npm run generate:books.
% canonical-source-hash: fnv1a64:7a8e745d630e52c5
% canonical-lessons: ZH-C289-ruguo, ZH-W289-guo-fruit, ZH-C289-ruguo-jiu, ZH-C289-would, ZH-C289-buran, ZH-W289-six-lines, ZH-R289-review-final-289

\chapter{If I Had the Time}
\label{ch:zh-if-i-had-the-time}
\hlchaptermodality{289}

\begin{chapteropening}
\textbf{By the end of this chapter:} \emph{I can say what I would do if things were different, and write six connected lines.}

Complete ZH-R289-review-final-289: one dialogue across the B1 core, after saying what would happen with \zh{如果} … \zh{就会} …, otherwise with \zh{不然}, and writing six connected lines to a friend about a trip that went wrong.
\end{chapteropening}

\section[\texorpdfstring{\zh{如果} (rúguǒ)}{如果 (rúguǒ)}]{\zh{如果} (rúguǒ) — if}

\label{lesson:ZH-C289-ruguo}



\begin{quote}
A new chapter, and it opens on the if of writing and careful speech.

\begin{itemize}
\raggedright
  \item \emph{Recall:} say \zh{要是下雨} \zh{我就不去了} — \emph{yàoshi xiàyǔ, wǒ jiù bú qù le} — \textbf{R1}, the chapter you have finished
  \item \emph{Recall:} say \zh{要是} — \emph{yàoshi}, if — \textbf{R2}
  \item \emph{Recall:} say \emph{jiéguǒ}, as a result — \textbf{R3}, durable
\end{itemize}
\end{quote}

\subsection*{You'll want to know: \zh{如果}}
\textbf{\zh{如果}} — \emph{rúguǒ} — \textquotedblleft{}if; in case\textquotedblright{}.

\begin{quote}
\textbf{\zh{如果你有时间} \zh{就来我家}} \emph{rúguǒ nǐ yǒu shíjiān, jiù lái wǒ jiā} — If you have time, come round.
\end{quote}

\zh{如果} is \zh{要是} in writing and in careful speech. The next lesson puts the second character in your hand; the first is the \zh{如} of \zh{比如}, and the second is the guǒ of \emph{jiéguǒ}.

\begin{cousinweb}
\textbf{\zh{如果}} (\emph{rúguǒ}) is built from \zh{如} and \zh{果}, both inherited from Middle Chinese, \emph{nyo} and \emph{kwaX} (Guangyun, in Baxter's transcription). The dictionary glosses the first with several words, among them 'to resemble', and the second 'resolute' and 'the fruit of a tree'. Whether 'if' was already among the first's senses the gloss does not make clear; 'if' for the pair is the modern word's sense.
\end{cousinweb}

\begin{grammarlens}[title={What you've built}]
\zh{如果} opens the condition as \zh{要是} does, and \zh{就} opens the result.

With a friend, \zh{要是} is more usual; in a note or a talk, \zh{如果}.
\end{grammarlens}

\subsection*{Your turn}
\begin{itemize}
\raggedright
  \item \emph{Say it:} \emph{rúguǒ}
  \item \emph{Say it:} \emph{rúguǒ}, once more
  \item \emph{Say it:} \zh{如果你有时间} \zh{就来我家} — \emph{rúguǒ nǐ yǒu shíjiān, jiù lái wǒ jiā} — If you have time, come round.
  \item \emph{Recall:} say \zh{要是}, then \zh{如果}: the everyday if, then the careful one
\end{itemize}

\begin{tcolorbox}[breakable,colback=teal!4,colframe=teal!35!black,title={Before you move on}]
What does \zh{如果} mean? (\textquotedblleft{}If\textquotedblright{}.) Which word means the same, more casually? (\textbf{\zh{要是}}.)
\end{tcolorbox}
\section[\texorpdfstring{\zh{果} (guǒ)}{果 (guǒ)}]{\zh{果} — \zh{日}, with \zh{木} running through it}

\label{lesson:ZH-W289-guo-fruit}



\begin{quote}
The word first, then the character inside it.

\begin{itemize}
\raggedright
  \item \emph{Recall:} say \emph{rúguǒ} — \textbf{R1}, one lesson back, before it decays
\end{itemize}

\begin{quote}
\zh{日} \zh{本}
\end{quote}

\emph{Read it:} each one aloud, with nothing said first — \emph{rì}, \emph{běn}
\end{quote}

\subsection*{Script — \zh{日}, with \zh{木} running through it}
\hlblockfigure{\detokenize{figures/ZH-W289-guo-fruit-filmstrip.pdf}}{How \zh{果} is written, stroke by stroke}

\begin{quote}
\zh{果}
\end{quote}

\textbf{Two pieces, and you write both already, one woven through the other.}

\textbf{\zh{日}, four strokes}, on top: the left side, the top that turns down the right side, the middle bar, the bottom bar. It is closed before anything else.

\textbf{\zh{木}, four strokes}, as in \zh{本}: the long horizontal under \zh{日}; then the vertical, which starts at the top of \zh{日} and runs down through it to the foot; then the left-falling stroke and the right-falling stroke.

4 + 4 = eight.

\subsection*{Your turn}
\begin{itemize}
\raggedright
  \item \emph{Write it:} \zh{果}, all eight, with the model in front of you
  \item \emph{Cover:} write it again from memory, and count the pieces, not the strokes
\end{itemize}

\begin{tcolorbox}[breakable,colback=teal!4,colframe=teal!35!black,title={Before you move on}]
How many strokes in \zh{果}? (\textbf{Eight}.) Which stroke runs through \zh{日}? (\textbf{The vertical of \zh{木}, the sixth}.) What is closed first? (\textbf{\zh{日}}.)

Source: \href{https://raw.githubusercontent.com/chanind/hanzi-writer-data/68d10a4b21150cae5e1ebbd223eed289cf32d90c/data/\%E6\%9E\%9C.json}{Hanzi Writer Data \zh{果}, pinned stroke paths and medians}.
\end{tcolorbox}
\section[\texorpdfstring{\zh{如果}…\zh{就}… (rúguǒ … jiù …)}{如果…就… (rúguǒ … jiù …)}]{\zh{如果}…\zh{就}… — if, then}

\label{lesson:ZH-C289-ruguo-jiu}



\begin{quote}
Last chapter's frame, with this chapter's word at its front.

\begin{itemize}
\raggedright
  \item \emph{Read it:} \zh{如果}, aloud — \emph{rúguǒ}
  \item \emph{Recall:} say \zh{要是你病了} \zh{就回家} — \emph{yàoshi nǐ bìng le, jiù huí jiā} — \textbf{R2}
  \item \emph{Recall:} say \zh{的话} — \emph{dehuà} — \textbf{R2}
\end{itemize}
\end{quote}

\begin{grammarlens}[title={Grammar Lens: \zh{如果} … \zh{就} …, if and then}]
\begin{quote}
\textbf{\zh{如果} opens the condition and \zh{就} opens the result: \zh{要是} … \zh{就} … for writing and careful speech.}
\end{quote}

\begin{quote}
\textbf{\zh{如果明天不下雨} \zh{我们就去山上}} \emph{rúguǒ míngtiān bú xiàyǔ, wǒmen jiù qù shān shàng} — If it doesn't rain tomorrow, we'll go up the mountain. \textbf{\zh{如果你有时间的话} \zh{就来我家}} \emph{rúguǒ nǐ yǒu shíjiān dehuà, jiù lái wǒ jiā} — If you have the time, come round.
\end{quote}

\zh{就} sits where it sat after \zh{要是}: after the subject, before the verb. \zh{的话} can close the condition here too.
\end{grammarlens}

\subsection*{Your turn}
Say these aloud:

\begin{itemize}
\raggedright
  \item \zh{如果明天不下雨} \zh{我们就去山上} — \emph{rúguǒ míngtiān bú xiàyǔ, wǒmen jiù qù shān shàng} — If it doesn't rain tomorrow, we'll go up the mountain.
  \item \zh{如果你有时间的话} \zh{就来我家} — \emph{rúguǒ nǐ yǒu shíjiān dehuà, jiù lái wǒ jiā} — If you have the time, come round.
  \item one plan for next week, with \zh{如果} … \zh{就} …
\end{itemize}

\begin{tcolorbox}[breakable,colback=teal!4,colframe=teal!35!black,title={Before you move on}]
Which word opens the result after \zh{如果}? (\textbf{\zh{就}}.) Where does it go? (\textbf{After the subject, before the verb}.)
\end{tcolorbox}
\section[\texorpdfstring{\zh{如果}…\zh{就会}… (rúguǒ … jiù huì …)}{如果…就会… (rúguǒ … jiù huì …)}]{\zh{如果}…\zh{就会}… — what would happen}

\label{lesson:ZH-C289-would}



\begin{quote}
A word from the A2 chapters on the future, after \zh{就}.

\begin{itemize}
\raggedright
  \item \emph{Read it:} \zh{如果明天不下雨} \zh{我们就去山上}, aloud — \emph{rúguǒ míngtiān bú xiàyǔ, wǒmen jiù qù shān shàng}
  \item \emph{Recall:} say \zh{明天会下雨} — \emph{míngtiān huì xiàyǔ} — \textbf{R3}, durable
  \item \emph{Recall:} say \zh{有钱} — \emph{yǒuqián}, rich — \textbf{R3}
\end{itemize}
\end{quote}

\begin{grammarlens}[title={Grammar Lens: \zh{如果} … \zh{就会} …, a result you picture}]
\begin{quote}
\textbf{\zh{会} after \zh{就} says what would happen: \zh{如果} … \zh{就会} … for a result you picture.}
\end{quote}

\begin{quote}
\textbf{\zh{如果我有钱} \zh{我就会去中国}} \emph{rúguǒ wǒ yǒu qián, wǒ jiù huì qù Zhōngguó} — If I had the money, I would go to China. \textbf{\zh{如果我有时间} \zh{我就会学开车}} \emph{rúguǒ wǒ yǒu shíjiān, wǒ jiù huì xué kāichē} — If I had the time, I would learn to drive.
\end{quote}

There is no new form here. It is the \zh{会} of \zh{明天会下雨}, a forecast, now forecasting a result you picture.

Nothing in the verb says whether the condition is real: \zh{如果我有钱} can be if I have the money, or if I had it. The rest of the conversation tells you.
\end{grammarlens}

\subsection*{Your turn}
Say these aloud:

\begin{itemize}
\raggedright
  \item \zh{如果我有钱} \zh{我就会去中国} — \emph{rúguǒ wǒ yǒu qián, wǒ jiù huì qù Zhōngguó} — If I had the money, I would go to China.
  \item \zh{如果我有时间} \zh{我就会学开车} — \emph{rúguǒ wǒ yǒu shíjiān, wǒ jiù huì xué kāichē} — If I had the time, I would learn to drive.
  \item what you would do if you had a free year, with \zh{如果} … \zh{就会} …
\end{itemize}

\begin{tcolorbox}[breakable,colback=teal!4,colframe=teal!35!black,title={Before you move on}]
What does \zh{会} add after \zh{就}? (\textbf{What would happen}.) Is it a new form? (\textbf{No: the \zh{会} of a forecast}.)
\end{tcolorbox}
\section[\texorpdfstring{\zh{不然} (bùrán)}{不然 (bùrán)}]{\zh{不然} (bùrán) — otherwise}

\label{lesson:ZH-C289-buran}



\begin{quote}
Before the new one: last lesson's picture, and a word for hurrying.

\begin{itemize}
\raggedright
  \item \emph{Read it:} \zh{如果我有钱} \zh{我就会去中国}, aloud — \emph{rúguǒ wǒ yǒu qián, wǒ jiù huì qù Zhōngguó}
  \item \emph{Recall:} say \zh{快} — \emph{kuài}, quick — \textbf{R3}, durable
\end{itemize}
\end{quote}

\subsection*{You'll want to know: \zh{不然}}
\textbf{\zh{不然}} — \emph{bùrán} — \textquotedblleft{}otherwise; or else\textquotedblright{}.

\begin{quote}
\textbf{\zh{快点儿} \zh{不然我们就晚了}} \emph{kuài diǎnr, bùrán wǒmen jiù wǎn le} — Hurry up, or we'll be late.
\end{quote}

\begin{quote}
\textbf{\zh{你得早点儿走} \zh{不然就晚了}} \emph{nǐ děi zǎo diǎnr zǒu, bùrán jiù wǎn le} — You have to leave early, or you'll be late.
\end{quote}

\zh{不} before the rising rán stays \emph{bù}. Every character in it is one you write.

\begin{cousinweb}
\textbf{\zh{不然}} (\emph{bùrán}) is built from \zh{不} 'not' and \zh{然} 'so', both inherited from Middle Chinese. The Guangyun lists \zh{不} under three readings, \emph{pjuw}, \emph{pjuwX} and \emph{pjut}, all glossed 'not'; bù is the regular modern sound of \emph{pjut}. \zh{然} is \emph{nyen}, glossed 'so, like this' (Guangyun, in Baxter's transcription). 'Not so', and so 'otherwise', is the plain sum of the two.
\end{cousinweb}

\begin{grammarlens}[title={What you've built}]
\zh{不然} comes after the advice and says what happens if it is not taken. \zh{就} often opens that result: \zh{不然我们就晚了}.
\end{grammarlens}

\subsection*{Your turn}
\begin{itemize}
\raggedright
  \item \emph{Say it:} \emph{bùrán}
  \item \emph{Say it:} \emph{bùrán}, once more
  \item \emph{Say it:} \zh{快点儿} \zh{不然我们就晚了} — \emph{kuài diǎnr, bùrán wǒmen jiù wǎn le} — Hurry up, or we'll be late.
  \item \emph{Say it:} \zh{你得早点儿走} \zh{不然就晚了} — \emph{nǐ děi zǎo diǎnr zǒu, bùrán jiù wǎn le} — You have to leave early, or you'll be late.
  \item \emph{Recall:} say \zh{然后}, then \zh{不然}: the same \zh{然}, two words
\end{itemize}

\begin{tcolorbox}[breakable,colback=teal!4,colframe=teal!35!black,title={Before you move on}]
What does \zh{不然} mean? (\textquotedblleft{}Otherwise\textquotedblright{}.) Where does it go? (\textbf{After the advice, before what happens if it is not taken}.)
\end{tcolorbox}
\section[(six lines)]{Six lines to a friend, a story with advice at the end}

\label{lesson:ZH-W289-six-lines}



\begin{quote}
A trip of yours went wrong: a train missed, a bag lost, rain where there should have been sun. A friend is about to make the same trip.

This lesson tells them what happened, and what to do.

\begin{itemize}
\raggedright
  \item \emph{Recall:} say \zh{不然} — \emph{bùrán}, otherwise — \textbf{R1}, one lesson back
  \item \emph{Recall:} say \zh{如果} … \zh{就会} … — \emph{rúguǒ … jiù huì …} — \textbf{R2}
\end{itemize}
\end{quote}

\subsection*{What to know first}
Six connected lines are not six sentences in a row: each leans on the one before, and the reader does not have to ask what happened next.

You own what does the leaning. \textbf{Order words}: \zh{先}, \zh{然后}, \zh{突然}, \zh{最后}. \textbf{The time at the front}: \zh{那天}, \zh{第二天}, \zh{的时候}. \textbf{The subject said once}: \zh{我} in the first line, then the order word or the time. \textbf{\zh{因为} … \zh{所以} …} gives the reason; \textbf{\zh{如果} … \zh{就} …} or \textbf{\zh{应该}} turns the story into advice.

Your reader is the friend. Your purpose is to save them the same trouble.

\subsection*{Writing — connected composition}
Write \textbf{six lines} to your friend. Here is a frame for each; the words are yours, and there is no model to copy.

\begin{itemize}
\raggedright
  \item Line 1: when and where, and what you did first, with \zh{先} and \zh{了}.
  \item Line 2: what came next, with \zh{然后}.
  \item Line 3: what went wrong, with \zh{突然}, or with \zh{的时候} and \zh{正在}.
  \item Line 4: how it ended, with \zh{最后}, or what you had never done before, with \zh{过}.
  \item Line 5: the reason it went wrong, with \zh{因为} … \zh{所以} ….
  \item Line 6: your advice, with \zh{如果} … \zh{就} … or \zh{应该}.
\end{itemize}

A word whose characters you have not written goes in pinyin, inside the line, with its tone marks. Every character comes from the 227 you have learned, each in the order you learned it, each in the same size of square.

No Chinese commas or full stops: leave a space between clauses, as this book prints them.

Write the six lines. Then read them once as your friend would, top to bottom: would they know what to do?

\begin{tcolorbox}[breakable,colback=teal!4,colframe=teal!35!black,title={Before you move on}]
Cover lines 1 and 2, and read line 3. Does it lean on them? (\textbf{It should: \zh{突然} breaks into what was going on}.)

Which line could your friend act on tomorrow? (\textbf{Line 6, the advice}.)
\end{tcolorbox}
\section[(dialogue)]{The B1 core — once more}

\label{lesson:ZH-R289-review-final-289}



\begin{quote}
The last lesson of the B1 core: one conversation that uses every chapter of it.

\begin{itemize}
\raggedright
  \item \emph{Recall:} say \emph{duōjiǔ}, how long — \textbf{R3}
  \item \emph{Recall:} say \emph{yíhàn}, a pity — \textbf{R3}
  \item \emph{Recall:} say \zh{不然} — \emph{bùrán}, otherwise — \textbf{R2}
\end{itemize}
\end{quote}

\subsection*{Your turn}
\begin{quote}
\textbf{A:} \zh{那天你做什么了}
\end{quote}

\begin{quote}
\emph{nà tiān nǐ zuò shénme le} — What did you do that day?
\end{quote}

\begin{quote}
\textbf{B:} \zh{我先去了书店} \zh{然后坐火车回家}
\end{quote}

\begin{quote}
\emph{wǒ xiān qù le shūdiàn, ránhòu zuò huǒchē huí jiā} — First the bookshop, then home by train.
\end{quote}

\begin{quote}
\textbf{B:} \zh{在火车上} \zh{我的手机突然不见了}
\end{quote}

\begin{quote}
\emph{zài huǒchē shàng, wǒ de shǒujī tūrán bú jiàn le} — On the train my phone suddenly went missing.
\end{quote}

\begin{quote}
\textbf{A:} \zh{你找回了吗}
\end{quote}

\begin{quote}
\emph{nǐ zhǎohuí le ma} — Did you get it back?
\end{quote}

\begin{quote}
\textbf{B:} \zh{最后找回了} \zh{我感觉很好}
\end{quote}

\begin{quote}
\emph{zuìhòu zhǎohuí le, wǒ gǎnjué hěn hǎo} — In the end I did. I feel fine.
\end{quote}

\begin{quote}
\textbf{B:} \emph{wǒ de zhèngjiàn yě bú jiàn le, wǒ děi mǎshàng guàshī} — My papers went missing too; I have to report them lost at once.
\end{quote}

\begin{quote}
\textbf{A:} \emph{wǒ juéde nǐ yīnggāi xiān liánxì wǒ} — I think you should have got in touch with me first.
\end{quote}

\begin{quote}
\textbf{B:} \emph{yīnwèi wǒ hěn lèi, suǒyǐ wàng le} — I was tired, so I forgot.
\end{quote}

\begin{quote}
\textbf{A:} \zh{以后你想当什么}
\end{quote}

\begin{quote}
\emph{yǐhòu nǐ xiǎng dāng shénme} — What do you want to be later on?
\end{quote}

\begin{quote}
\textbf{B:} \emph{wǒ de mèngxiǎng shì dāng yīshēng, wǒ de mùbiāo shì shuō hǎo Hànyǔ} — My dream is to be a doctor; my goal is to speak good Chinese.
\end{quote}

\begin{quote}
\textbf{A:} \emph{nǐ xué le duōjiǔ} — How long have you studied?
\end{quote}

\begin{quote}
\textbf{B:} \emph{liǎng nián, wǒ méi qù guo Zhōngguó, tài yíhàn le} — Two years. I've never been to China, sadly.
\end{quote}

\begin{quote}
\textbf{A:} \zh{明天去山上吗} \zh{去山上怎么走}
\end{quote}

\begin{quote}
\emph{míngtiān qù shān shàng ma, qù shān shàng zěnme zǒu} — Going up the mountain tomorrow? How do you get there?
\end{quote}

\begin{quote}
\textbf{B:} \zh{要看天气} \zh{要是下雨的话} \zh{我就不去了}
\end{quote}

\begin{quote}
\emph{yào kàn tiānqì, yàoshi xiàyǔ dehuà, wǒ jiù bú qù le} — It depends on the weather. If it rains, I won't go.
\end{quote}

\begin{quote}
\textbf{A:} \zh{如果不下雨} \zh{我们就坐火车去}
\end{quote}

\begin{quote}
\emph{rúguǒ bú xiàyǔ, wǒmen jiù zuò huǒchē qù} — If it doesn't rain, we'll go by train.
\end{quote}

\begin{quote}
\textbf{A:} \zh{快点儿} \zh{不然我们就晚了}
\end{quote}

\begin{quote}
\emph{kuài diǎnr, bùrán wǒmen jiù wǎn le} — Hurry, or we'll be late.
\end{quote}

\begin{itemize}
\raggedright
  \item \emph{Read it:} the lines printed in characters aloud, then the rest
  \item \emph{Say it:} the whole dialogue, both parts
  \item \emph{Say it:} B's lines, with your own answers
\end{itemize}

\begin{tcolorbox}[breakable,colback=teal!4,colframe=teal!35!black,title={Before you move on}]
First, then, last? (\textbf{\zh{先}, \zh{然后}, \zh{最后}}.) Should? (\textbf{\zh{应该}}.) Gone missing? (\textbf{\zh{不见了}}.) How long? (\emph{duōjiǔ}.) If, then? (\textbf{\zh{要是} … \zh{就} …, \zh{如果} … \zh{就} …}.) Would? (\textbf{\zh{就会}}.) Otherwise? (\textbf{\zh{不然}}.)
\end{tcolorbox}
